\chapter{Differential forms, connections, and gauge transformations}\label{ch:forms}

\section{Forms, wedge, exterior derivative}\label{sec:forms-basics}

The reader of \cite{Schwartz2014QFT} is already fluent in components and index contractions: for a tensor $T_{\mu\nu}$ the operation $\partial_{[\mu}T_{\nu\rho]}$ produces another tensor of definite transformation law, Gauss's theorem pairs a divergence with a surface integral, and the chiral anomaly of Chapter~30 is written as $\varepsilon^{\mu\nu\rho\sigma}\tr F_{\mu\nu}F_{\rho\sigma}$. The language of differential forms is a bookkeeping system that absorbs exactly this kind of antisymmetrization and orientation data into the notation, so that instead of chasing factorials and epsilon tensors one manipulates forms directly. Nothing new is going on physically. What changes is that identities like the Bianchi identity or the statement $\tr(F\wedge F) = d\omega_{\CS}$ become one-line calculations rather than index gymnastics.

\medskip

\noindent The material of this section is entirely standard and can be read off any textbook in differential geometry; Nakahara \cite{Nakahara2003GTP} is the pedagogical reference for a physics audience. We give definitions here in enough detail to fix conventions for the rest of the paper and to settle every sign. The key technical result, collected in Lemma~\ref{lem:ch2-leibniz-nilpotent} below, is the pair of identities $d(\alpha\wedge\beta) = (d\alpha)\wedge\beta + (-1)^p\alpha\wedge(d\beta)$ and $d^2 = 0$. Later chapters use these identities constantly.

\subsection{Local definition of a $p$-form}

Fix a smooth $n$-manifold $M$. In a local chart with coordinates $x^\mu$, a \emph{differential $p$-form} is an expression
\begin{equation}\label{eq:form-local}
\omega \;=\; \frac{1}{p!}\,\omega_{\mu_1\cdots\mu_p}(x)\,dx^{\mu_1}\wedge\cdots\wedge dx^{\mu_p},
\end{equation}
in which the components $\omega_{\mu_1\cdots\mu_p}$ are smooth functions on the chart, totally antisymmetric in their $p$ indices, and the symbols $dx^\mu$ are the basis one-forms dual to $\partial/\partial x^\mu$. The wedge product $\wedge$ will be defined momentarily; for the moment it is enough to record that it is associative, bilinear, and satisfies $dx^\mu\wedge dx^\nu = -dx^\nu \wedge dx^\mu$.

\medskip

\noindent The space of smooth $p$-forms on $M$ is denoted $\Omega^p(M)$. By convention, $\Omega^0(M) = C^\infty(M)$ is the space of smooth functions on $M$; these are the ``$0$-forms''. For $p > n$ the top form-degree on an $n$-manifold is exceeded, and $\Omega^p(M) = 0$ because any antisymmetric tensor with more slots than there are coordinates must vanish by the pigeonhole principle applied to the indices.

\begin{remark}[On the factor $1/p!$]\label{rem:pfactorial}
The $1/p!$ prefactor in \eqref{eq:form-local} makes the sum over all $p$-tuples $(\mu_1,\ldots,\mu_p)$ count each unordered set of indices exactly once after accounting for the $p!$ ways of permuting the $dx^\mu$'s. Without this prefactor, an expression like $T_{\mu\nu}\,dx^\mu\wedge dx^\nu$ with $T_{\mu\nu} = -T_{\nu\mu}$ antisymmetric would double-count because the pairs $(\mu,\nu)$ and $(\nu,\mu)$ contribute the same quantity. Some authors absorb the $1/p!$ by summing only over ordered multi-indices $\mu_1 < \mu_2 < \cdots < \mu_p$; we use the Einstein-summation convention so the factorials appear explicitly.
\end{remark}

\subsection{The wedge product}

The wedge product $\wedge : \Omega^p(M) \times \Omega^q(M) \to \Omega^{p+q}(M)$ is defined on basis one-forms by imposing associativity and antisymmetry $dx^\mu \wedge dx^\nu = -dx^\nu \wedge dx^\mu$, extended to higher-degree basis forms by juxtaposition
\begin{equation*}
(dx^{\mu_1}\wedge\cdots\wedge dx^{\mu_p}) \wedge (dx^{\nu_1}\wedge\cdots\wedge dx^{\nu_q}) \;=\; dx^{\mu_1}\wedge\cdots\wedge dx^{\mu_p}\wedge dx^{\nu_1}\wedge\cdots\wedge dx^{\nu_q},
\end{equation*}
and extended bilinearly to arbitrary smooth forms. The wedge of two general forms $\alpha \in \Omega^p(M)$ and $\beta \in \Omega^q(M)$ is then
\begin{equation}\label{eq:wedge-components}
(\alpha \wedge \beta)_{\mu_1\cdots\mu_{p+q}} \;=\; \frac{(p+q)!}{p!\,q!}\,\alpha_{[\mu_1\cdots\mu_p}\,\beta_{\mu_{p+1}\cdots\mu_{p+q}]},
\end{equation}
where the square brackets denote total antisymmetrization over the enclosed indices with weight $1$ (so the antisymmetrizer is $\tfrac{1}{k!}$ times the signed sum over permutations, where $k$ is the number of antisymmetrized indices). The reader who prefers component gymnastics can take \eqref{eq:wedge-components} as the definition; we will not need it in this form after this subsection.

\begin{proposition}[Graded commutativity]\label{prop:graded-comm}
For $\alpha \in \Omega^p(M)$ and $\beta \in \Omega^q(M)$,
\begin{equation}\label{eq:ch2-graded-comm}
\alpha \wedge \beta \;=\; (-1)^{pq}\,\beta \wedge \alpha.
\end{equation}
\end{proposition}

\begin{proof}
The component functions $\alpha_{\mu_1\cdots\mu_p}$ and $\beta_{\nu_1\cdots\nu_q}$ commute as scalars. The sign comes entirely from reordering the wedge of basis one-forms:
\begin{equation*}
dx^{\mu_1}\wedge\cdots\wedge dx^{\mu_p}\wedge dx^{\nu_1}\wedge\cdots\wedge dx^{\nu_q} \;=\; (-1)^{pq}\,dx^{\nu_1}\wedge\cdots\wedge dx^{\nu_q}\wedge dx^{\mu_1}\wedge\cdots\wedge dx^{\mu_p}.
\end{equation*}
To see the sign, move the leftmost $dx^{\mu_1}$ past all $q$ basis one-forms of $\beta$: this is $q$ adjacent transpositions, each contributing a factor $-1$, for a total of $(-1)^q$. Repeat for $dx^{\mu_2}, \dots, dx^{\mu_p}$: each of the remaining $p-1$ one-forms moves past all $q$ of $\beta$'s basis one-forms and contributes $(-1)^q$. The total sign is $\left((-1)^q\right)^p = (-1)^{pq}$.
\end{proof}

\begin{remark}[Two useful specializations]\label{rem:wedge-specials}
If $\alpha$ is an odd-degree form and we take $\beta = \alpha$ in \eqref{eq:ch2-graded-comm}, then $\alpha \wedge \alpha = -\alpha \wedge \alpha$, so $\alpha \wedge \alpha = 0$. In particular, the square of any one-form on $M$ is zero. This is why a single basis one-form $dx^\mu$ satisfies $dx^\mu \wedge dx^\mu = 0$ and why the wedge of a one-form $A$ with itself vanishes in the \emph{scalar-valued} case. Lie-algebra-valued one-forms behave differently, and we return to this in Section~\ref{sec:lie-forms} below.

On the other hand, even-degree forms commute: if $\alpha$ has degree $p = 2k$, then $\alpha \wedge \beta = \beta \wedge \alpha$ for any $\beta$.
\end{remark}

\subsection{The exterior derivative}

The exterior derivative $d : \Omega^p(M) \to \Omega^{p+1}(M)$ is defined in local coordinates by
\begin{equation}\label{eq:d-def}
d\omega \;=\; \frac{1}{p!}\,\partial_\nu \omega_{\mu_1\cdots\mu_p}\,dx^\nu \wedge dx^{\mu_1}\wedge\cdots\wedge dx^{\mu_p}.
\end{equation}
On a 0-form, that is, a smooth function $f$, this specializes to $df = \partial_\nu f\, dx^\nu$, the familiar object. On a one-form $A = A_\mu\,dx^\mu$ it gives
\begin{equation*}
dA \;=\; \partial_\nu A_\mu\, dx^\nu \wedge dx^\mu \;=\; \tfrac{1}{2}(\partial_\mu A_\nu - \partial_\nu A_\mu)\,dx^\mu \wedge dx^\nu,
\end{equation*}
where the second equality antisymmetrizes by adding and subtracting; explicitly,
\begin{align*}
\partial_\nu A_\mu\, dx^\nu\wedge dx^\mu
&= \tfrac{1}{2}(\partial_\nu A_\mu - \partial_\mu A_\nu)\,dx^\nu\wedge dx^\mu + \tfrac{1}{2}(\partial_\nu A_\mu + \partial_\mu A_\nu)\,dx^\nu\wedge dx^\mu \\
&= \tfrac{1}{2}(\partial_\nu A_\mu - \partial_\mu A_\nu)\,dx^\nu\wedge dx^\mu,
\end{align*}
the symmetric piece killing the antisymmetric wedge. Renaming $\nu\leftrightarrow\mu$ in the surviving piece yields $dA = \tfrac{1}{2}(\partial_\mu A_\nu - \partial_\nu A_\mu)\,dx^\mu\wedge dx^\nu$, as claimed. The reader will recognize this as the Maxwell field strength in the form-valued notation we shall exploit throughout.

\medskip

\noindent Note that \eqref{eq:d-def} as written is a statement about a coordinate patch; the non-trivial fact, which we do not prove here because it would detain us from the physics, is that this local formula extends to a globally well-defined operator $d$ on $\Omega^p(M)$. The proof consists in checking that the right-hand side of \eqref{eq:d-def} transforms correctly under a change of chart, which it does because the basis one-forms $dx^\mu$ transform as a cotangent-vector does and the partial derivative in the first slot is paired with a new one-form $dx^\nu$ in a way that cancels the chart-transition non-tensorial term. This is covered in Nakahara \cite{Nakahara2003GTP}.

\begin{example}[Exterior derivative of a 2-form on $\R^3$]\label{ex:d-2form}
Let $\omega = P\,dy\wedge dz + Q\,dz\wedge dx + R\,dx\wedge dy$ on $\R^3$, where $P, Q, R$ are smooth functions of $(x,y,z)$. Compute
\begin{align*}
d\omega
&= dP \wedge dy\wedge dz + dQ\wedge dz \wedge dx + dR\wedge dx\wedge dy \\
&= \left(\partial_x P\,dx + \partial_y P\,dy + \partial_z P\,dz\right)\wedge dy\wedge dz + \cdots \\
&= \partial_x P\,dx\wedge dy\wedge dz + \partial_y Q\,dy\wedge dz\wedge dx + \partial_z R\,dz\wedge dx\wedge dy \\
&= (\partial_x P + \partial_y Q + \partial_z R)\,dx\wedge dy\wedge dz,
\end{align*}
where in the third line the terms $\partial_y P\,dy\wedge dy\wedge dz = 0$ and $\partial_z P\,dz \wedge dy \wedge dz = 0$ have dropped out because $dx^\mu\wedge dx^\mu = 0$, and similarly for $Q$ and $R$; and in the fourth line the basis three-form $dy\wedge dz\wedge dx$ has been rewritten as $dx\wedge dy\wedge dz$ using two adjacent transpositions (sign $+1$), and $dz\wedge dx\wedge dy$ likewise. The result reproduces the divergence of the vector field $(P,Q,R)$: the identity $d\omega = (\nabla \cdot \vec{V})\,\mathrm{dvol}$ for $\omega$ the ``flux'' 2-form of $\vec{V}$ is an example of the general pattern that $d$ in forms language packages the differential operators grad, curl, and div.
\end{example}

\subsection{The Leibniz rule and the nilpotency of $d$}

Two identities involving $d$ are used on every page of the rest of this paper.

\begin{lemma}[Leibniz rule and nilpotency]\label{lem:ch2-leibniz-nilpotent}
For $\alpha \in \Omega^p(M)$ and $\beta \in \Omega^q(M)$,
\begin{align}
d(\alpha \wedge \beta) &\;=\; (d\alpha)\wedge \beta + (-1)^p\,\alpha \wedge (d\beta), \label{eq:ch2-leibniz}\\[1mm]
d^2 &\;=\; 0. \label{eq:ch2-nilpotent}
\end{align}
\end{lemma}

\begin{proof}
Both identities follow from the definition \eqref{eq:d-def} by direct computation. We work locally in a chart; the results are coordinate-invariant because $d$ is.

\medskip

\noindent \emph{The Leibniz rule \eqref{eq:ch2-leibniz}.} Write $\alpha = \frac{1}{p!}\alpha_{\mu_1\cdots\mu_p}\,dx^{\mu_1}\wedge\cdots\wedge dx^{\mu_p}$ and $\beta = \frac{1}{q!}\beta_{\nu_1\cdots\nu_q}\,dx^{\nu_1}\wedge\cdots\wedge dx^{\nu_q}$. Their wedge has components
\begin{equation*}
\alpha \wedge \beta \;=\; \frac{1}{p!\,q!}\,\alpha_{\mu_1\cdots\mu_p}\beta_{\nu_1\cdots\nu_q}\,dx^{\mu_1}\wedge\cdots\wedge dx^{\mu_p}\wedge dx^{\nu_1}\wedge\cdots\wedge dx^{\nu_q}.
\end{equation*}
Apply $d$ according to the definition \eqref{eq:d-def}, which places a $dx^\rho$ at the left and a $\partial_\rho$ acting on the component:
\begin{equation*}
d(\alpha \wedge \beta) \;=\; \frac{1}{p!\,q!}\,\partial_\rho\!\left(\alpha_{\mu_1\cdots\mu_p}\beta_{\nu_1\cdots\nu_q}\right) dx^\rho \wedge dx^{\mu_1}\wedge\cdots\wedge dx^{\mu_p}\wedge dx^{\nu_1}\wedge\cdots\wedge dx^{\nu_q}.
\end{equation*}
The ordinary product rule splits the derivative into two pieces. The piece in which $\partial_\rho$ acts on $\alpha$ reads
\begin{equation*}
\frac{1}{p!\,q!}\,\partial_\rho \alpha_{\mu_1\cdots\mu_p}\,\beta_{\nu_1\cdots\nu_q}\,dx^\rho \wedge dx^{\mu_1}\wedge\cdots\wedge dx^{\mu_p}\wedge dx^{\nu_1}\wedge\cdots\wedge dx^{\nu_q},
\end{equation*}
which is precisely $(d\alpha)\wedge\beta$ by the definitions of $d\alpha$ and the wedge.

\medskip

\noindent The piece in which $\partial_\rho$ acts on $\beta$ gives
\begin{equation*}
\frac{1}{p!\,q!}\,\alpha_{\mu_1\cdots\mu_p}\,\partial_\rho\beta_{\nu_1\cdots\nu_q}\,dx^\rho \wedge dx^{\mu_1}\wedge\cdots\wedge dx^{\mu_p}\wedge dx^{\nu_1}\wedge\cdots\wedge dx^{\nu_q}.
\end{equation*}
To recognize this as $(-1)^p\,\alpha\wedge(d\beta)$, we must move the leftmost $dx^\rho$ past the $p$ basis one-forms $dx^{\mu_i}$ to reach the position immediately to the left of $dx^{\nu_1}$. Each transposition of $dx^\rho$ past one $dx^{\mu_i}$ costs a factor $-1$; moving past all $p$ of them costs $(-1)^p$. After the move, the wedge has the shape $dx^{\mu_1}\wedge\cdots\wedge dx^{\mu_p}\wedge dx^\rho\wedge dx^{\nu_1}\wedge\cdots\wedge dx^{\nu_q}$, and the full expression reads
\begin{equation*}
(-1)^p\,\frac{1}{p!\,q!}\,\alpha_{\mu_1\cdots\mu_p}\,\partial_\rho\beta_{\nu_1\cdots\nu_q}\,dx^{\mu_1}\wedge\cdots\wedge dx^{\mu_p}\wedge dx^\rho \wedge dx^{\nu_1}\wedge\cdots\wedge dx^{\nu_q},
\end{equation*}
which is exactly $(-1)^p\,\alpha\wedge(d\beta)$ by the definitions of $d\beta$ and the wedge. Summing the two pieces proves \eqref{eq:ch2-leibniz}.

\medskip

\noindent \emph{Nilpotency \eqref{eq:ch2-nilpotent}.} Apply $d$ twice to $\omega \in \Omega^p(M)$:
\begin{equation*}
d^2\omega \;=\; \frac{1}{p!}\,\partial_\sigma\partial_\rho \omega_{\mu_1\cdots\mu_p}\,dx^\sigma \wedge dx^\rho \wedge dx^{\mu_1}\wedge\cdots\wedge dx^{\mu_p}.
\end{equation*}
Mixed partials commute for smooth functions, so $\partial_\sigma\partial_\rho\,\omega_{\mu_1\cdots\mu_p}$ is symmetric under $\rho\leftrightarrow\sigma$. On the other hand, the basis two-form $dx^\sigma\wedge dx^\rho$ is antisymmetric under the same swap. Writing the sum explicitly over the two contributions $(\sigma,\rho) = (a,b)$ and $(\sigma,\rho) = (b,a)$ for each unordered pair $a \ne b$,
\begin{align*}
\partial_a\partial_b\omega_{\cdots}\,dx^a\wedge dx^b + \partial_b\partial_a\omega_{\cdots}\,dx^b\wedge dx^a
&= \partial_a\partial_b\omega_{\cdots}\,dx^a\wedge dx^b - \partial_a\partial_b\omega_{\cdots}\,dx^a\wedge dx^b \\
&= 0.
\end{align*}
The diagonal contributions $(\sigma,\rho) = (a,a)$ vanish because $dx^a \wedge dx^a = 0$ by Remark~\ref{rem:wedge-specials}. Summing over all pairs gives $d^2\omega = 0$.
\end{proof}

\begin{remark}[On the symmetry-versus-antisymmetry argument]\label{rem:sym-anti}
The nilpotency proof uses the general principle that the contraction of a symmetric object with an antisymmetric object in the same two indices is zero. Concretely, if $S_{\rho\sigma} = S_{\sigma\rho}$ and $A^{\rho\sigma} = -A^{\sigma\rho}$, then
\begin{equation*}
S_{\rho\sigma}A^{\rho\sigma} \;=\; S_{\sigma\rho}A^{\rho\sigma} \;=\; -S_{\sigma\rho}A^{\sigma\rho} \;=\; -S_{\rho\sigma}A^{\rho\sigma},
\end{equation*}
where the first equality uses the symmetry of $S$, the second uses the antisymmetry of $A$, and the third renames dummy indices. Therefore $S_{\rho\sigma}A^{\rho\sigma} = -S_{\rho\sigma}A^{\rho\sigma}$, and $S_{\rho\sigma}A^{\rho\sigma} = 0$. This little identity is the engine of $d^2 = 0$ and recurs throughout the form calculus.
\end{remark}

\subsection{Maxwell theory in forms language}\label{sec:forms-maxwell}

To give the machinery a concrete face, we rewrite elementary electrodynamics in forms language and read off two familiar facts as immediate consequences of $d^2 = 0$.

\begin{example}[Maxwell theory]\label{ex:maxwell-forms}
The four-vector potential $A_\mu(x)$ of \cite[Ch.~8]{Schwartz2014QFT} is the component of a globally defined one-form $A = A_\mu\,dx^\mu$ on spacetime. The field strength tensor $F_{\mu\nu} = \partial_\mu A_\nu - \partial_\nu A_\mu$ is the component of the two-form
\begin{equation}\label{eq:F-as-dA}
F \;=\; dA \;=\; \tfrac{1}{2}F_{\mu\nu}\,dx^\mu\wedge dx^\nu,
\end{equation}
where the second equality was computed in the paragraph following \eqref{eq:d-def}. The usual decomposition of $F_{\mu\nu}$ into electric and magnetic components $E_i = F_{0i}$, $B_i = \tfrac{1}{2}\varepsilon_{ijk}F^{jk}$ is just the choice of a time coordinate and a spatial basis; as a differential form, $F$ needs no such choice.

\medskip

\noindent Apply $d$ to both sides of \eqref{eq:F-as-dA}. By Lemma~\ref{lem:ch2-leibniz-nilpotent},
\begin{equation}\label{eq:bianchi-forms}
dF \;=\; d(dA) \;=\; d^2 A \;=\; 0.
\end{equation}
This is the \emph{Bianchi identity} in forms language. In components, \eqref{eq:bianchi-forms} expands to
\begin{equation*}
0 \;=\; dF \;=\; \tfrac{1}{2}\partial_\rho F_{\mu\nu}\,dx^\rho\wedge dx^\mu\wedge dx^\nu,
\end{equation*}
which forces $\partial_{[\rho}F_{\mu\nu]} = 0$, equivalently the two homogeneous Maxwell equations $\nabla\cdot\vec{B} = 0$ and $\partial_t\vec{B} + \nabla\times\vec{E} = 0$. Crucially, \eqref{eq:bianchi-forms} is true for \emph{any} $A$, whether or not the inhomogeneous Maxwell equations are satisfied.

\medskip

\noindent The inhomogeneous equations $\partial^\mu F_{\mu\nu} = J_\nu$, by contrast, are genuine equations of motion and appear in forms language as
\begin{equation}\label{eq:maxwell-inhom}
d \star F \;=\; \star J,
\end{equation}
where $\star : \Omega^p(M) \to \Omega^{n-p}(M)$ is the Hodge star and $J = J_\mu\, dx^\mu$ the current one-form. The Hodge star is not defined by $d$ and $\wedge$ alone: its construction requires a choice of metric on $M$, since the definition involves raising indices. Consequently \eqref{eq:maxwell-inhom} is metric-dependent and its content changes if we deform the metric, whereas \eqref{eq:bianchi-forms} does not.

\medskip

\noindent This is the pattern that will organize the rest of the paper: \emph{identities involving only $d$ and $\wedge$ are topological and metric-free, while equations involving $\star$ are dynamical and metric-dependent.} The Chern--Simons action of Chapter~\ref{ch:chern-simons} uses only $d$ and $\wedge$ and is therefore a topological field theory; the Yang--Mills action $\int \tr(F \wedge \star F)$ uses the Hodge star and is metric-dependent, which is what makes it dynamical.
\end{example}

\begin{example}[Gauge invariance of $F$]\label{ex:gauge-invariance-F}
A $U(1)$ gauge transformation of the potential $A_\mu \to A_\mu + \partial_\mu \lambda$ reads in forms language as
\begin{equation*}
A \;\longmapsto\; A + d\lambda,
\end{equation*}
where $\lambda$ is a smooth function, i.e.\ a $0$-form. The field strength transforms as
\begin{equation*}
F \;\longmapsto\; d(A + d\lambda) \;=\; dA + d^2\lambda \;=\; dA \;=\; F,
\end{equation*}
using $d^2 = 0$ once. So $F$ is gauge-invariant: this is a one-line consequence of nilpotency.
\end{example}

\medskip

\noindent This completes the algebraic foundation. In Section~\ref{sec:integration-stokes} we develop integration and Stokes's theorem, which is the other half of the toolkit. De~Rham cohomology, which the pair of identities in Lemma~\ref{lem:ch2-leibniz-nilpotent} makes possible, is defined there.


\section{Integration, Stokes's theorem, and de~Rham cohomology}\label{sec:integration-stokes}

The exterior derivative of Section~\ref{sec:forms-basics} is a differential operator; pairing it with a matching integration theory produces the technical statement we use in the rest of the paper, namely that the integral of an exact form over a closed manifold vanishes. That one identity is the engine behind the gauge-invariance and level-quantization arguments of Chapter~\ref{ch:chern-simons}; it is also the reason de~Rham cohomology carries topological information.

\subsection{Orientation and integration of top forms}

Before integrating we must orient. A smooth $n$-manifold $M$ is \emph{orientable} if it admits an atlas whose coordinate-change Jacobians all have positive determinant; a choice of such an atlas is an \emph{orientation}. The orientable manifolds one meets in the rest of this paper --- spheres, tori, three-manifolds with boundary in physical spacetime, and surfaces dressed with a spin structure --- are all orientable, and we fix an orientation once and for all on each.

\medskip

\noindent Let $M$ be an oriented $n$-manifold and $\omega \in \Omega^n(M)$ a top form. In a single oriented chart with coordinates $x^1, \ldots, x^n$, one writes
\begin{equation}\label{eq:top-form-local}
\omega = f(x)\,dx^1 \wedge dx^2 \wedge \cdots \wedge dx^n,
\end{equation}
and the integral of $\omega$ over the image of the chart is
\begin{equation}\label{eq:top-form-integral-local}
\int_U \omega = \int_{\phi(U)} f(x)\,d^n x,
\end{equation}
where $\phi: U \to \phi(U) \subseteq \R^n$ is the chart map and the right-hand side is an ordinary Lebesgue integral. The sign convention is forced by the orientation: if one changes chart to coordinates $y^1, \ldots, y^n$ with Jacobian $J = \det(\partial y / \partial x)$, the form has components $f(x) = \tilde f(y(x)) \cdot |J|\, \mathrm{sgn}(J)$, and the positive-Jacobian hypothesis on the chart atlas makes $\mathrm{sgn}(J) = +1$. The integral is therefore coordinate-independent.

\medskip

\noindent On a manifold covered by more than one chart, the integral is defined by partition of unity: pick a partition $\{\rho_i\}_i$ subordinate to the atlas, so that $\sum_i \rho_i = 1$ with each $\rho_i$ supported in a single chart, and set
\begin{equation}\label{eq:pou-integral}
\int_M \omega := \sum_i \int_{U_i} \rho_i\, \omega.
\end{equation}
Standard arguments (see \cite{Nakahara2003GTP}) show that the right-hand side of \eqref{eq:pou-integral} is independent of the chosen partition and of the chosen atlas within its orientation class. We will not redo these arguments here; the reader unfamiliar with partitions of unity can take \eqref{eq:pou-integral} as a black-box definition. What matters for the physics downstream is that we have a well-defined real number $\int_M \omega$ assigned to every pair $(M, \omega)$ with $M$ a compact oriented $n$-manifold and $\omega$ a top form on it.

\medskip

\noindent Integration is linear: for top forms $\omega, \omega'$ and scalars $a, b \in \R$,
\begin{equation*}
\int_M (a\omega + b\omega') = a \int_M \omega + b \int_M \omega'.
\end{equation*}
We also need integration of a $p$-form over a $p$-dimensional oriented submanifold $\Sigma \subseteq M$: the integral is defined by pulling back $\omega$ to $\Sigma$ along the inclusion $\iota : \Sigma \hookrightarrow M$ and then integrating the top form $\iota^* \omega \in \Omega^p(\Sigma)$ via \eqref{eq:pou-integral}. In physics notation this is the same as saying $\int_\Sigma \omega = \int_\Sigma \omega|_\Sigma$, restricting $\omega$ to $\Sigma$ using the embedding.

\begin{remark}[Why orientation matters]\label{rem:orientation}
Reversing the orientation of $M$ flips the sign of every integral:
\begin{equation*}
\int_{M^{\mathrm{op}}} \omega = - \int_M \omega.
\end{equation*}
This is where the parity-odd character of the Chern--Simons action of Chapter~\ref{ch:chern-simons} comes from: the action $S_{\CS}[A] = \tfrac{k}{4\pi}\int_M \omega_{\CS}(A)$ flips sign under $M \mapsto M^{\mathrm{op}}$. The quantum weight $e^{i S_{\CS}}$ is therefore not invariant under orientation reversal, which is the technical reason Chern--Simons theory distinguishes left from right.
\end{remark}

\subsection{Stokes's theorem}

The central analytic input of this section is Stokes's theorem, which relates the integral of an exterior derivative over a manifold to the integral of the form itself over the boundary.

\begin{theorem}[Stokes]\label{thm:ch2-stokes}
Let $M$ be a compact oriented $n$-manifold with boundary $\partial M$, equipped with the boundary orientation induced from $M$, and let $\omega \in \Omega^{n-1}(M)$. Then
\begin{equation}\label{eq:ch2-stokes}
\int_M d\omega = \int_{\partial M} \omega.
\end{equation}
If $M$ is closed, so that $\partial M = \emptyset$, then
\begin{equation}\label{eq:ch2-stokes-closed}
\int_M d\omega = 0.
\end{equation}
\end{theorem}

A proof is given in Nakahara \cite{Nakahara2003GTP}, along with the construction of the induced orientation on the boundary. The argument proceeds in three steps: reduce to the case where $\omega$ is supported in a single chart by a partition of unity; in that chart, reduce to $M = [0, 1]^n$ or to a half-space; finally, on the half-space, reduce to the fundamental theorem of calculus applied one coordinate at a time. We will not repeat the argument here.

\medskip

\noindent We shall rely on the closed-manifold corollary \eqref{eq:ch2-stokes-closed} constantly. Three immediate consequences are worth spelling out.

\begin{corollary}[Exact forms integrate to zero on closed manifolds]\label{cor:exact-to-zero}
If $M$ is a compact oriented $n$-manifold without boundary and $\omega = d\eta \in \Omega^n(M)$ is exact, then
\begin{equation*}
\int_M \omega = \int_M d\eta = 0.
\end{equation*}
\end{corollary}

\begin{proof}
Immediate from \eqref{eq:ch2-stokes-closed}.
\end{proof}

\begin{corollary}[Homology dependence]\label{cor:homology-dependence}
Let $\omega \in \Omega^p(M)$ be closed, and let $\Sigma, \Sigma' \subseteq M$ be two compact oriented $p$-submanifolds without boundary that are the boundary of a common $(p+1)$-chain, so that $\Sigma - \Sigma' = \partial \Omega$ for some $(p+1)$-chain $\Omega$. Then
\begin{equation*}
\int_\Sigma \omega = \int_{\Sigma'} \omega.
\end{equation*}
\end{corollary}

\begin{proof}
By hypothesis and Theorem~\ref{thm:ch2-stokes},
\begin{equation*}
\int_\Sigma \omega - \int_{\Sigma'} \omega = \int_{\partial \Omega} \omega = \int_\Omega d\omega = 0,
\end{equation*}
the last equality because $\omega$ is closed.
\end{proof}

\noindent Corollary~\ref{cor:homology-dependence} is the reason the integer in the level-quantization calculation of Chapter~\ref{ch:chern-simons} depends only on homology: the gauge-transformed Chern--Simons three-form differs from the original by the sum of an exact form and a topological Wess--Zumino piece, so its integral over a closed three-manifold is determined by the homology class of the manifold.

\begin{corollary}[Homotopy invariance for closed forms]\label{cor:homotopy-invariance}
Let $N$ be a compact oriented $p$-manifold without boundary, $M$ any smooth manifold, $\omega \in \Omega^p(M)$ closed, and $f_0, f_1 : N \to M$ smoothly homotopic. Then
\begin{equation*}
\int_N f_0^* \omega = \int_N f_1^* \omega.
\end{equation*}
\end{corollary}

A proof uses Cartan's magic formula $\mathcal{L}_X = \iota_X d + d\iota_X$ to write the difference as an exact form on $N$ and then applies Stokes; see Nakahara \cite{Nakahara2003GTP}. We will not reproduce the argument.

\medskip

\noindent Corollary~\ref{cor:homotopy-invariance} is the reason ``the integral of a normalized closed form is a homotopy invariant'' is a theorem rather than a slogan, and it is the reason the degree of a map between oriented manifolds of the same dimension is well defined.

\subsection{Worked example: the winding number of $f : S^1 \to S^1$}

We unpack a single example in full, because the level-quantization calculation of Chapter~\ref{ch:chern-simons} is structurally identical to it.

\begin{example}[Winding number]\label{ex:ch2-winding}
Let $S^1 = \R / 2\pi\Z$ be the circle with angular coordinate $\theta$. The one-form $d\theta$ is \emph{globally well-defined} on $S^1$ even though $\theta$ itself is multi-valued: in a chart where $\theta$ ranges over $(a, a + 2\pi)$ for some $a$, the expression $d\theta$ is the differential of the local coordinate, and two choices of chart differ by $\theta \mapsto \theta + 2\pi k$, under which $d(\theta + 2\pi k) = d\theta + 0 = d\theta$. So $d\theta$ is globally defined. It is also closed, $d(d\theta) = d^2\theta = 0$, by Lemma~\ref{lem:ch2-leibniz-nilpotent}.

\medskip

\noindent Let $f : S^1 \to S^1$ be any smooth map. Consider
\begin{equation}\label{eq:winding-int}
n(f) := \frac{1}{2\pi} \int_{S^1} f^* d\theta.
\end{equation}
We claim that $n(f) \in \Z$ and that $n(f)$ depends only on the homotopy class of $f$.

\medskip

\noindent For the integrality, parametrize the source circle by $t \in [0, 2\pi]$ with endpoints identified. In this chart the function $f$ lifts to a smooth real-valued function $\tilde f : [0, 2\pi] \to \R$ satisfying $\tilde f(2\pi) - \tilde f(0) = 2\pi\, n$ for some integer $n$, because $f$ takes $t = 0$ and $t = 2\pi$ to the same point of the target circle. The pullback is
\begin{equation*}
f^* d\theta = d(\tilde f) = \tilde f'(t)\, dt,
\end{equation*}
and the integral \eqref{eq:winding-int} becomes
\begin{equation*}
n(f) = \frac{1}{2\pi} \int_0^{2\pi} \tilde f'(t)\,dt = \frac{1}{2\pi} \left(\tilde f(2\pi) - \tilde f(0)\right) = \frac{1}{2\pi} \cdot 2\pi n = n \in \Z.
\end{equation*}
The evaluation uses the fundamental theorem of calculus for the single-variable integral in the middle and the lifting property of $\tilde f$ at the end.

\medskip

\noindent For the homotopy invariance, let $f_s : S^1 \to S^1$ be a smooth one-parameter family with $s \in [0, 1]$. Corollary~\ref{cor:homotopy-invariance} applies because $S^1$ is compact and oriented without boundary and $d\theta$ is closed: the integrals at $s = 0$ and $s = 1$ are equal.

\medskip

\noindent The integer $n(f)$ is the \emph{winding number} or \emph{degree} of $f$. It can be computed explicitly for the map $f(t) = e^{int}$, which lifts to $\tilde f(t) = nt$ and winds $n$ times around the target circle.

\medskip

\noindent \emph{Forward pointer.} Example~\ref{ex:ch2-winding} is the prototype for every topological quantization argument in this paper. A differential form --- here $d\theta$ --- is globally well-defined and closed, though it is locally the derivative of a function that is not globally well-defined; its integral over a closed manifold produces an integer. In Chapter~\ref{ch:chern-simons} the form will be $\tr\left((g^{-1}dg)^{\wedge 3}\right)$ on the group $SU(2) \cong S^3$, and the integer will be the winding number of a gauge transformation $g : M^3 \to SU(2)$; the logical structure is identical to the one here, with $S^1$ replaced by $SU(2)$ and $\theta$ replaced by the Maurer--Cartan form of \S\ref{sec:gauge-transf}. The integer is what forces the Chern--Simons level $k$ to be an integer, and it is the reason Witten could compute link invariants in $S^3$ from a path integral.
\end{example}

\subsection{Closed and exact forms; de~Rham cohomology}\label{sec:deRham}

We can now package the algebraic facts of Section~\ref{sec:forms-basics} and the analytic facts of this section into a single invariant of the manifold. A $p$-form $\omega$ is \emph{closed} if $d\omega = 0$ and \emph{exact} if there exists a $(p-1)$-form $\eta$ with $\omega = d\eta$. By Lemma~\ref{lem:ch2-leibniz-nilpotent}, every exact form is closed, because $d(d\eta) = d^2\eta = 0$.

\begin{definition}[De~Rham cohomology]\label{def:deRham}
The $p$-th \emph{de Rham cohomology} of a smooth manifold $M$ is the quotient
\begin{equation}\label{eq:deRham-def}
H^p_{\mathrm{dR}}(M) := \frac{\kernel(d : \Omega^p(M) \to \Omega^{p+1}(M))}{\image(d : \Omega^{p-1}(M) \to \Omega^p(M))}.
\end{equation}
A class $[\omega] \in H^p_{\mathrm{dR}}(M)$ is represented by a closed $p$-form $\omega$, and two closed forms represent the same class if and only if they differ by an exact form.
\end{definition}

\noindent The quotient in \eqref{eq:deRham-def} is well defined because $\image\, d \subseteq \kernel\, d$ by nilpotency of $d$. Dimensions $0, 1, \ldots, n$ on an $n$-manifold give a sequence of real vector spaces $H^0_{\mathrm{dR}}(M), H^1_{\mathrm{dR}}(M), \ldots, H^n_{\mathrm{dR}}(M)$, known collectively as the de~Rham cohomology of $M$. On a compact manifold, each $H^p_{\mathrm{dR}}(M)$ is finite-dimensional and the sum of their dimensions is a topological invariant (the Euler characteristic, up to signs), although we will not use the latter fact.

\begin{theorem}[de~Rham; cited]\label{thm:deRham-iso}
For any smooth manifold $M$, the de~Rham cohomology $H^p_{\mathrm{dR}}(M)$ is canonically isomorphic to the singular cohomology $H^p(M; \R)$ with real coefficients.
\end{theorem}

The isomorphism sends a closed form $\omega$ to the functional $[\Sigma] \mapsto \int_\Sigma \omega$ on singular $p$-cycles. The non-trivial content of the theorem is that this functional depends only on the de~Rham class of $\omega$ (which follows from Corollary~\ref{cor:homology-dependence} above) and that the resulting map on cohomology classes is an isomorphism. A proof is in Nakahara \cite{Nakahara2003GTP}.

\medskip

\noindent A second cohomological fact we cite without proof is the Poincar\'e lemma: on a contractible open set, every closed form of positive degree is exact.

\begin{proposition}[Poincar\'e lemma; cited]\label{prop:poincare-lemma}
Let $U \subseteq \R^n$ be an open, star-shaped set, and let $\omega \in \Omega^p(U)$ with $p \geq 1$ be closed. Then there exists $\eta \in \Omega^{p-1}(U)$ such that $\omega = d\eta$.
\end{proposition}

An explicit homotopy operator $h : \Omega^p(U) \to \Omega^{p-1}(U)$ satisfying $dh + hd = \mathrm{id}$ on positive-degree forms provides the proof, which is given in Nakahara \cite{Nakahara2003GTP}. The Poincar\'e lemma is the reason every topological statement in this paper that involves an obstruction has global input: locally, closed forms are exact, so obstructions to exactness are global topological features of the manifold.

\begin{remark}[Why de~Rham cohomology drives the physics]\label{rem:deRham-physics}
Several of the integer-valued invariants in this paper are, in disguise, pairings of de~Rham classes with homology classes. The second Chern number $\tfrac{1}{8\pi^2}\int_M \tr(F \wedge F)$ of a gauge bundle on a closed 4-manifold $M$ is the pairing of the de~Rham class $[\tr(F \wedge F)]$ with the fundamental class of $M$; its integrality is a cohomological statement. The Chern--Simons level $k$ is integer-quantized because the gauge-transformation shift of $S_{\CS}$ is expressible as the integral of a closed form that represents an integer cohomology class. The obstruction to gauging a higher-form symmetry discussed in Chapter~\ref{ch:gensym} lives in an integer cohomology group; the de~Rham isomorphism is what lets us write it as the integral of a differential form. This is why we set up de~Rham cohomology in this primer and not as an isolated appendix.
\end{remark}

\medskip

\noindent This concludes the integration and cohomology material. Section~\ref{sec:lie-forms} adds the Lie-algebra-valued extension we need for non-abelian gauge theory, and Section~\ref{sec:gauge-transf} closes the chapter by deriving the transformation law of the field strength. After that, Chapter~\ref{ch:axiomatic} returns to the topological side: the Atiyah--Segal axiomatic framework builds on top of the calculus we have just developed.


\section{Lie-algebra-valued forms and the non-abelian field strength}\label{sec:lie-forms}

Chapters~\ref{ch:chern-simons}, \ref{ch:bf}, \ref{ch:dw}, \ref{ch:gensym}, and \ref{ch:anyons} all work with connections valued in a non-abelian Lie algebra. Two features make this case different from the scalar-valued one of the previous sections: the wedge product no longer annihilates an odd-degree form against itself, and a convention choice is forced on us by the interplay between physicists' Hermitian generators and geometers' anti-Hermitian generators. We settle both here once and for all.

\subsection{Conventions: anti-Hermitian generators}\label{sec:anti-hermitian-conv}

Let $G$ be a compact Lie group with Lie algebra $\mathfrak{g}$. In Schwartz \cite{Schwartz2014QFT}, the generators $T^a$ of $\mathfrak{g}$ are \emph{Hermitian} matrices in the fundamental representation, satisfying
\begin{equation}\label{eq:schwartz-generators}
[T^a, T^b] \;=\; i f^{abc} T^c, \qquad \tr(T^a T^b) \;=\; \tfrac{1}{2} \delta^{ab},
\end{equation}
with $f^{abc}$ the real structure constants. Throughout the rest of this paper we use a different convention, standard in the differential-geometry and TQFT literature \cite{Freed1995CSnotes, Witten1989Jones}: the \emph{anti-Hermitian} basis
\begin{equation}\label{eq:anti-herm-def}
t^a \;:=\; -i\, T^a, \qquad t^{a\dagger} \;=\; -t^a.
\end{equation}

The translation between the two conventions is mechanical but must be kept explicit, because the same physical connection $A$ appears in the paper in both forms depending on which chapter the reader started from.

\begin{proposition}[Commutation relations and trace in the anti-Hermitian basis]\label{prop:anti-herm-identities}
The anti-Hermitian generators $t^a = -iT^a$ satisfy
\begin{align}
[t^a, t^b] &\;=\; f^{abc}\, t^c, \label{eq:anti-herm-bracket}\\[1mm]
\tr(t^a t^b) &\;=\; -\tfrac{1}{2}\,\delta^{ab}. \label{eq:anti-herm-trace}
\end{align}
The absence of the $i$ in \eqref{eq:anti-herm-bracket}, compared to the $i f^{abc}$ appearing on the Hermitian side, is the reason the anti-Hermitian convention is natural on the geometry side: the structure constants $f^{abc}$ appear with no extraneous factor of $i$.
\end{proposition}

\begin{proof}
For \eqref{eq:anti-herm-bracket}, using \eqref{eq:anti-herm-def},
\begin{equation*}
[t^a, t^b] \;=\; [(-i)T^a,\, (-i)T^b] \;=\; (-i)^2 [T^a, T^b] \;=\; -\, i f^{abc} T^c,
\end{equation*}
by \eqref{eq:schwartz-generators}. Now use $T^c = it^c$, which follows from $t^c = -iT^c$ on inverting:
\begin{equation*}
-\,i f^{abc} T^c \;=\; -\,i f^{abc}\,(i t^c) \;=\; -\,i^2 f^{abc}\, t^c \;=\; f^{abc}\, t^c.
\end{equation*}

For \eqref{eq:anti-herm-trace},
\begin{equation*}
\tr(t^a t^b) \;=\; \tr\!\left((-iT^a)(-iT^b)\right) \;=\; (-i)^2 \tr(T^a T^b) \;=\; -\,\tr(T^a T^b) \;=\; -\,\tfrac{1}{2}\delta^{ab},
\end{equation*}
by \eqref{eq:schwartz-generators}.
\end{proof}

\begin{definition}[Connection one-form]\label{def:connection-one-form}
A \emph{connection one-form} on $M$ (in the anti-Hermitian convention) is a Lie-algebra-valued one-form
\begin{equation}\label{eq:A-components}
A \;=\; A_\mu^a\, t^a\, dx^\mu \;=\; A_\mu\, dx^\mu, \qquad A_\mu \,:=\, A_\mu^a\, t^a \,\in\, \mathfrak{g},
\end{equation}
with smooth real components $A_\mu^a(x)$. The translation back to Schwartz's Hermitian convention is
\begin{equation}\label{eq:A-schwartz-translation}
A_{\mathrm{here}} \;=\; -\,i\, A_{\mathrm{Schwartz}}.
\end{equation}
\end{definition}

\begin{remark}[Box: the translation rule]\label{rem:translation-box}
For the remainder of this paper we use the anti-Hermitian convention exclusively. To translate any of our formulas to Schwartz's Hermitian convention, substitute
\begin{equation*}
A \;\longrightarrow\; -\,i\,A, \qquad F \;\longrightarrow\; -\,i\,F, \qquad t^a \;\longrightarrow\; -\,i\,T^a.
\end{equation*}
The field strength translation $F_{\mathrm{here}} = -\,i\,F_{\mathrm{Schwartz}}$ will be verified explicitly at the end of this section; see the computation after \eqref{eq:F-nonabelian-components}.
\end{remark}

\subsection{Wedge products of Lie-algebra-valued forms}

For Lie-algebra-valued $p$- and $q$-forms, the wedge product is defined by taking the scalar wedge on the form side and matrix multiplication on the Lie-algebra side:
\begin{equation}\label{eq:lie-wedge-def}
\alpha \wedge \beta \;:=\; \alpha^a \wedge \beta^b\, t^a t^b,
\end{equation}
where $\alpha = \alpha^a\, t^a$ and $\beta = \beta^b\, t^b$ with $\alpha^a, \beta^b$ scalar forms of degrees $p$ and $q$ respectively. The right-hand side of \eqref{eq:lie-wedge-def} is a $(p+q)$-form valued in $\mathfrak{g} \otimes \mathfrak{g}$; the matrix product $t^a t^b$ is taken in whatever representation of $\mathfrak{g}$ is in use, and we suppress the representation index throughout because we will rarely need it.

\medskip

\noindent Because matrix multiplication is noncommutative, the wedge product of Lie-algebra-valued forms no longer satisfies the graded-commutativity property \eqref{eq:ch2-graded-comm} on the nose. Instead, the relation is
\begin{equation}\label{eq:lie-graded-comm}
\alpha \wedge \beta \;=\; (-1)^{pq}\, \beta^b \wedge \alpha^a\, t^a t^b \;=\; (-1)^{pq}\,(\beta \wedge \alpha)\bigr|_{t^a t^b \leftrightarrow t^b t^a},
\end{equation}
where the last factor indicates that swapping $\alpha$ and $\beta$ on the form side forces us to also reverse the order of the generator product, leaving a commutator to account for. This motivates the following.

\begin{definition}[Graded commutator of Lie-algebra-valued forms]\label{def:graded-commutator}
For Lie-algebra-valued forms $\alpha$ and $\beta$ of degrees $p$ and $q$ respectively,
\begin{equation}\label{eq:graded-commutator-def}
[\alpha, \beta] \;:=\; \alpha \wedge \beta \,-\, (-1)^{pq}\, \beta \wedge \alpha.
\end{equation}
\end{definition}

\noindent The graded commutator is well defined because the two terms on the right-hand side are Lie-algebra-valued forms of the same degree $p+q$.

\begin{proposition}[$A\wedge A$ in terms of the commutator]\label{prop:AwedgeA}
For a connection one-form $A = A_\mu^a\, t^a\, dx^\mu$,
\begin{equation}\label{eq:AwedgeA}
A \wedge A \;=\; \tfrac{1}{2}\,[A_\mu, A_\nu]\, dx^\mu \wedge dx^\nu \;=\; \tfrac{1}{2}\, f^{abc}\, A_\mu^a\, A_\nu^b\, t^c\, dx^\mu \wedge dx^\nu.
\end{equation}
In particular, $A \wedge A \neq 0$ for a non-abelian Lie algebra, in contrast with the scalar case of Remark~\ref{rem:wedge-specials}.
\end{proposition}

\begin{proof}
By the definition \eqref{eq:lie-wedge-def},
\begin{equation*}
A \wedge A \;=\; A_\mu^a\, A_\nu^b\, t^a t^b\, dx^\mu \wedge dx^\nu.
\end{equation*}
The strategy is the standard one: the coefficient $A_\mu^a A_\nu^b t^a t^b$ has a symmetric and an antisymmetric part under the swap $(\mu, a) \leftrightarrow (\nu, b)$, and only the antisymmetric part survives pairing with the antisymmetric wedge $dx^\mu \wedge dx^\nu$. Write
\begin{equation*}
A_\mu^a A_\nu^b\, t^a t^b \;=\; \tfrac{1}{2}\,A_\mu^a A_\nu^b\, \left(t^a t^b + t^b t^a\right) \,+\, \tfrac{1}{2}\,A_\mu^a A_\nu^b\, \left(t^a t^b - t^b t^a\right).
\end{equation*}
The first term is symmetric under $(\mu, a)\leftrightarrow(\nu, b)$ because $A_\mu^a A_\nu^b \{t^a, t^b\}$ is unchanged by simultaneously relabeling $(\mu, a)\leftrightarrow(\nu, b)$ (the anticommutator $\{t^a, t^b\} = t^at^b + t^bt^a$ is symmetric in $a\leftrightarrow b$, and the scalar coefficients commute). Paired against $dx^\mu \wedge dx^\nu$, which is antisymmetric under $\mu \leftrightarrow \nu$, it produces zero by the symmetric-antisymmetric-contraction principle of Remark~\ref{rem:sym-anti}.

\medskip

\noindent The surviving antisymmetric piece is
\begin{equation*}
\tfrac{1}{2}\,A_\mu^a A_\nu^b\,[t^a, t^b]\, dx^\mu \wedge dx^\nu.
\end{equation*}
Using \eqref{eq:anti-herm-bracket}, $[t^a, t^b] = f^{abc}\, t^c$, so this equals
\begin{equation*}
\tfrac{1}{2}\,f^{abc}\,A_\mu^a\, A_\nu^b\, t^c\, dx^\mu \wedge dx^\nu,
\end{equation*}
which is the right-hand side of \eqref{eq:AwedgeA}. To collapse the form into the $[A_\mu, A_\nu]$ notation, observe that
\begin{equation*}
[A_\mu, A_\nu] \;=\; [A_\mu^a t^a,\, A_\nu^b t^b] \;=\; A_\mu^a A_\nu^b\, [t^a, t^b] \;=\; f^{abc}\, A_\mu^a A_\nu^b\, t^c,
\end{equation*}
where the second equality uses the bilinearity of the commutator in the scalar coefficients. Substituting back gives the middle expression of \eqref{eq:AwedgeA}.
\end{proof}

\begin{remark}[The factor of two between $[A,A]$ and $A\wedge A$]\label{rem:AwedgeA-factor}
Applying the graded commutator of Definition~\ref{def:graded-commutator} to $\alpha = \beta = A$ (both one-forms, so $p = q = 1$) gives
\begin{equation*}
[A, A] \;=\; A \wedge A - (-1)^{1\cdot 1}\, A \wedge A \;=\; A\wedge A + A\wedge A \;=\; 2\, A \wedge A.
\end{equation*}
This is a source of confusion in the literature: if one sees $[A, A]$ in a paper and interprets it via the commutator-of-one-forms definition, one gets twice what one would get using $A \wedge A$ directly. We always write $A \wedge A$ rather than $\tfrac{1}{2}[A, A]$ to avoid the ambiguity.
\end{remark}

\begin{remark}[The abelian case]\label{rem:abelian-case}
For $G = U(1)$ the algebra is one-dimensional with $f^{abc} = 0$, and \eqref{eq:AwedgeA} specializes to $A \wedge A = 0$. The abelian curvature is simply $F = dA$, exactly as in Maxwell theory (Example~\ref{ex:maxwell-forms}); the nonabelian cubic correction drops out. Chapter~\ref{ch:chern-simons} uses this fact when treating the abelian Chern--Simons theory of Laughlin states: the cubic term $\tfrac{2}{3}A\wedge A\wedge A$ in the Chern--Simons density vanishes for $U(1)$, and the action reduces to $\frac{k}{4\pi}\int A\wedge dA$.
\end{remark}

\subsection{The non-abelian field strength}

With the conventions fixed, the non-abelian field strength takes a particularly clean form.

\begin{definition}[Non-abelian field strength]\label{def:F-nonabelian}
The \emph{field strength} of a connection $A$ is the Lie-algebra-valued two-form
\begin{equation}\label{eq:F-nonabelian-def}
F \;:=\; dA + A \wedge A.
\end{equation}
\end{definition}

\begin{proposition}[Components of the field strength]\label{prop:F-components}
In components,
\begin{equation}\label{eq:F-nonabelian-components}
F \;=\; \tfrac{1}{2}\,F_{\mu\nu}\, dx^\mu \wedge dx^\nu,
\qquad
F_{\mu\nu} \;=\; \partial_\mu A_\nu - \partial_\nu A_\mu + [A_\mu, A_\nu].
\end{equation}
\end{proposition}

\begin{proof}
The exterior derivative of $A = A_\mu^a\, t^a\, dx^\mu$ is, by \eqref{eq:d-def},
\begin{equation*}
dA \;=\; \partial_\nu A_\mu^a\, t^a\, dx^\nu \wedge dx^\mu.
\end{equation*}
The coefficient $\partial_\nu A_\mu^a\, t^a$ is Lie-algebra-valued with no symmetry in $(\mu, \nu)$; the same symmetric-antisymmetric decomposition as in the abelian case gives
\begin{equation*}
dA \;=\; \tfrac{1}{2}\,\bigl(\partial_\mu A_\nu^a - \partial_\nu A_\mu^a\bigr)\, t^a\, dx^\mu \wedge dx^\nu \;=\; \tfrac{1}{2}\,\bigl(\partial_\mu A_\nu - \partial_\nu A_\mu\bigr)\, dx^\mu \wedge dx^\nu,
\end{equation*}
where the second equality absorbs the generator index into $A_\mu = A_\mu^a t^a$. Adding to this the cubic piece $A\wedge A = \tfrac{1}{2}[A_\mu, A_\nu]\, dx^\mu \wedge dx^\nu$ computed in Proposition~\ref{prop:AwedgeA} produces
\begin{equation*}
F \;=\; dA + A\wedge A \;=\; \tfrac{1}{2}\,\bigl(\partial_\mu A_\nu - \partial_\nu A_\mu + [A_\mu, A_\nu]\bigr)\, dx^\mu \wedge dx^\nu,
\end{equation*}
which is \eqref{eq:F-nonabelian-components}.
\end{proof}

\begin{remark}[Comparison with Schwartz]\label{rem:schwartz-comparison}
Schwartz's non-abelian field strength, in Hermitian conventions, reads
\begin{equation}\label{eq:F-schwartz}
F_{\mu\nu}^{\mathrm{Sch}} \;=\; \partial_\mu A_\nu^{\mathrm{Sch}} - \partial_\nu A_\mu^{\mathrm{Sch}} \,-\, i\, [A_\mu^{\mathrm{Sch}}, A_\nu^{\mathrm{Sch}}].
\end{equation}
This differs from \eqref{eq:F-nonabelian-components} by the factor of $-i$ in front of the commutator. We verify that our convention \eqref{eq:F-nonabelian-components} and Schwartz's \eqref{eq:F-schwartz} are related by the translation \eqref{eq:A-schwartz-translation}. Substituting $A_\mu = -i\, A_\mu^{\mathrm{Sch}}$ into $F_{\mu\nu} = \partial_\mu A_\nu - \partial_\nu A_\mu + [A_\mu, A_\nu]$:
\begin{align*}
\partial_\mu A_\nu - \partial_\nu A_\mu &\;=\; -\,i\,\bigl(\partial_\mu A_\nu^{\mathrm{Sch}} - \partial_\nu A_\mu^{\mathrm{Sch}}\bigr), \\[1mm]
[A_\mu, A_\nu] &\;=\; \bigl[-i A_\mu^{\mathrm{Sch}},\, -i A_\nu^{\mathrm{Sch}}\bigr] \;=\; (-i)^2\,[A_\mu^{\mathrm{Sch}}, A_\nu^{\mathrm{Sch}}] \;=\; -\,[A_\mu^{\mathrm{Sch}}, A_\nu^{\mathrm{Sch}}].
\end{align*}
Summing,
\begin{align*}
F_{\mu\nu} &\;=\; -\,i\,\bigl(\partial_\mu A_\nu^{\mathrm{Sch}} - \partial_\nu A_\mu^{\mathrm{Sch}}\bigr) \,-\, [A_\mu^{\mathrm{Sch}}, A_\nu^{\mathrm{Sch}}] \\[1mm]
&\;=\; -\,i\,\Bigl(\partial_\mu A_\nu^{\mathrm{Sch}} - \partial_\nu A_\mu^{\mathrm{Sch}} \,-\, i\,[A_\mu^{\mathrm{Sch}}, A_\nu^{\mathrm{Sch}}]\Bigr) \\[1mm]
&\;=\; -\,i\,F_{\mu\nu}^{\mathrm{Sch}},
\end{align*}
where in the second equality we factored out the overall $-i$ and used $-1 = -i \cdot i$ on the commutator piece. This confirms the translation $F_{\mathrm{here}} = -i\,F_{\mathrm{Schwartz}}$ promised in Remark~\ref{rem:translation-box}.
\end{remark}

\begin{example}[Abelian check]\label{ex:abelian-F-check}
For $G = U(1)$, Remark~\ref{rem:abelian-case} gives $A \wedge A = 0$, so $F = dA$ and $F_{\mu\nu} = \partial_\mu A_\nu - \partial_\nu A_\mu$. This matches both Schwartz's Maxwell formula (after the trivial translation $A \to -iA$, which in the abelian $U(1)$ case amounts to a multiplicative rescaling by $-i$) and Example~\ref{ex:maxwell-forms} of this chapter. The two conventions agree on the \emph{shape} of $F$ when the algebra is one-dimensional; they disagree on overall factors of $-i$, which is harmless in the abelian case because physical quantities like $\sigma_{xy} = ke^2/h$ depend only on the level $k$ and not on whether $A$ carries an $i$.
\end{example}

\begin{remark}[Abelian $\tr$]\label{rem:abelian-trace}
For the abelian $U(1)$ gauge field, there is only one generator (the single-element $\mathfrak{u}(1) = i\R$ algebra), and the symbol $\tr$ is nothing more than a multiplicative scalar fixed by normalization. We adopt the convention $\tr \to 1$ for $U(1)$ calculations, which makes $\tr(F \wedge F) = F \wedge F$ and $\tr(A \wedge dA) = A \wedge dA$ with no prefactors. This is the convention used in Chapter~\ref{ch:chern-simons} for the abelian Laughlin action.
\end{remark}

\medskip

\noindent The non-abelian field strength satisfies a Bianchi identity
\begin{equation}\label{eq:nonabelian-bianchi}
dF + A \wedge F - F \wedge A \;=\; 0,
\end{equation}
equivalently $DF := dF + [A, F] = 0$ for the graded commutator of Definition~\ref{def:graded-commutator}. We will not use \eqref{eq:nonabelian-bianchi} in the rest of the paper, so we record it for completeness and move on. Section~\ref{sec:principal-bundles} records the minimum principal-bundle vocabulary needed for Chapter~\ref{ch:chern-simons} and for the discussion of nontrivial bundles in Chapter~\ref{ch:gensym}.


\section{Principal bundles, connections, and global issues}\label{sec:principal-bundles}

The objects we have been calling ``connection one-forms'' in Section~\ref{sec:lie-forms} have a more intrinsic home: they are local expressions of a global object called a \emph{connection} on a \emph{principal bundle} over the base manifold $M$. We will not develop the theory of principal bundles here. The minimum we need for Chapter~\ref{ch:chern-simons} is (i) a clear statement that, under the hypotheses of that chapter, the bundle is trivial and one can work globally with $A \in \Omega^1(M, \mathfrak{g})$; and (ii) a hook for the discussion of nontrivial bundles in Chapter~\ref{ch:gensym}, where magnetic monopoles and instantons are bundle-theoretic in origin. A full treatment is in Nakahara \cite{Nakahara2003GTP}.

\subsection{Informal definition}\label{sec:principal-bundle-def}

A \emph{principal $G$-bundle} over an $n$-manifold $M$ is a smooth manifold $P$ equipped with a free right action of the Lie group $G$ and a smooth projection $\pi : P \to M$, subject to the condition that $M$ is the quotient $P/G$ and that $P$ is \emph{locally trivial}: around each point $x \in M$ there exists an open neighborhood $U$ and a diffeomorphism $\phi_U : \pi^{-1}(U) \xrightarrow{\sim} U \times G$ commuting with both $\pi$ (after projection to the first factor) and the right $G$-action. A \emph{local section} over $U$ is a smooth map $s : U \to P$ with $\pi \circ s = \mathrm{id}_U$, and the trivialization $\phi_U$ determines such a section by $s(x) := \phi_U^{-1}(x, e)$.

The picture to hold in mind: at each point of $M$ there is a copy of $G$ (the fiber $\pi^{-1}(x) \cong G$), and the local trivializations glue these copies together by a ``transition function''. The transitions between two local trivializations $\phi_U, \phi_V$ over overlapping patches take the form $g_{UV} : U \cap V \to G$, acting on the fiber by right multiplication. A principal $G$-bundle is \emph{trivial} if it admits a single global trivialization, equivalently if there is a global section $s : M \to P$.

\subsection{Connection one-form on the total space}\label{sec:connection-on-total-space}

A \emph{connection} on $P$ is a $\mathfrak{g}$-valued one-form $\omega \in \Omega^1(P, \mathfrak{g})$ satisfying two axioms. First, at each $p \in P$, the \emph{vertical condition}: for every $\xi \in \mathfrak{g}$ with associated fundamental vector field $\xi^\# \in TP$ (the infinitesimal generator of the right $G$-action),
\begin{equation}\label{eq:vertical-axiom}
\omega_p(\xi^\#_p) \;=\; \xi.
\end{equation}
Second, \emph{equivariance} under the right action of $g \in G$, with $R_g : P \to P$ denoting the action,
\begin{equation}\label{eq:equivariance-axiom}
R_g^* \,\omega \;=\; \mathrm{Ad}(g^{-1})\, \omega.
\end{equation}

These two conditions are non-negotiable: any $\omega$ satisfying them is called a connection, and any connection satisfies them. We will not reproduce their motivation here; Nakahara \cite{Nakahara2003GTP} develops the story in full. All we need downstream is that they exist and that their local expressions are the $A = A_\mu^a t^a\, dx^\mu$ objects we have been using.

\begin{definition}[Local connection via a section]\label{def:local-connection}
Given a local section $s : U \to P$ and a connection $\omega$ on $P$, the \emph{local connection} on $U$ is the pullback
\begin{equation}\label{eq:local-connection-def}
A \;:=\; s^* \omega \;\in\; \Omega^1(U, \mathfrak{g}).
\end{equation}
Writing $A = A_\mu^a t^a\, dx^\mu$ in a coordinate chart on $U$ recovers the component notation of Definition~\ref{def:connection-one-form} in Section~\ref{sec:lie-forms}.
\end{definition}

\noindent The local connection $A$ depends on the section $s$. If we pick a different section $s' = s \cdot g$ (related by a $G$-valued function $g : U \to G$), the equivariance axiom \eqref{eq:equivariance-axiom} produces a specific transformation law for $A$: this is exactly the gauge transformation of Section~\ref{sec:gauge-transf}. So gauge transformations are, intrinsically, changes of local section.

\subsection{Curvature on the total space and on the base}\label{sec:curvature-on-P}

The \emph{curvature} of a connection $\omega$ on $P$ is the two-form
\begin{equation}\label{eq:total-space-curvature}
\Omega \;:=\; d\omega + \omega \wedge \omega \;\in\; \Omega^2(P, \mathfrak{g}).
\end{equation}
This is the global version of the formula $F = dA + A \wedge A$ of Section~\ref{sec:lie-forms}. A basic calculation (Nakahara \cite{Nakahara2003GTP}) shows that $\Omega$ is \emph{horizontal} in the sense that $\Omega_p(\xi^\#_p, \cdot) = 0$ for every $\xi \in \mathfrak{g}$, and \emph{equivariant} under the right action. These two properties together imply that $\Omega$ descends to a well-defined 2-form on the base $M$, provided one has chosen a local section, and the descended object is precisely
\begin{equation}\label{eq:F-as-pullback}
F \;=\; s^* \Omega \;=\; dA + A \wedge A,
\end{equation}
the field strength of Proposition~\ref{prop:F-components}. The equality \eqref{eq:F-as-pullback} connects the total-space formulation we have just sketched to the local formulation we already use.

\subsection{When is the bundle trivial?}

We now state the triviality fact that Chapter~\ref{ch:chern-simons} uses as a standing assumption.

\begin{proposition}[Trivial bundles over 3-manifolds; cited]\label{prop:trivial-bundle-3-manifold}
Let $M$ be a closed oriented smooth three-manifold, and let $G$ be a compact simple simply connected Lie group (e.g., $G = SU(N)$). Then every principal $G$-bundle over $M$ is trivial.
\end{proposition}

The precise classification statement is that principal $G$-bundles over $M$ are in bijection with $[M, BG]$, the set of homotopy classes of maps into the classifying space $BG$. Obstruction theory plus the homotopy groups $\pi_k(BG) = \pi_{k-1}(G)$ for $k \geq 1$ reduce Proposition~\ref{prop:trivial-bundle-3-manifold} to the facts $\pi_0(G) = \pi_1(G) = 0$ (which come from simply connectedness) and $\pi_2(G) = 0$ (a classical theorem: no compact simple Lie group has $\pi_2 \neq 0$). Nakahara \cite{Nakahara2003GTP} gives the argument in full. For our purposes it is enough that, in Chapter~\ref{ch:chern-simons}, the principal $SU(N)$-bundle over any closed oriented $M^3$ is trivial and we can work globally with $A \in \Omega^1(M, \mathfrak{g})$.

\begin{remark}[What fails in Proposition~\ref{prop:trivial-bundle-3-manifold}]\label{rem:when-triviality-fails}
Dropping either the dimensional hypothesis or the simply-connectedness of $G$ can produce nontrivial bundles. Two canonical examples will matter in later chapters.
\begin{enumerate}
\item \emph{$U(1)$-bundles over $S^2$.} The group $G = U(1)$ is not simply connected; $\pi_1(U(1)) = \Z$. Principal $U(1)$-bundles over $S^2$ are classified by $\pi_1(U(1)) = \Z$, and the integer class of a bundle is the first Chern number $\frac{1}{2\pi}\int_{S^2} F$. This integer is the magnetic monopole charge of Dirac; we return to it in Chapter~\ref{ch:gensym}, where gauging a 1-form symmetry of electromagnetism realizes it as a charged 0-form symmetry object.
\item \emph{$SU(N)$-bundles over $S^4$.} Over a 4-dimensional base, the classification uses $\pi_3(G)$; for $G$ compact simple simply connected one has $\pi_3(G) = \Z$, and the integer class is the instanton number $\frac{1}{8\pi^2}\int_{S^4} \tr(F \wedge F)$. This is the second Chern number whose integrality underwrites the chiral-anomaly calculation of \cite[Ch.~30]{Schwartz2014QFT} and, re-told in differential-form language, the level quantization of Chapter~\ref{ch:chern-simons}.
\end{enumerate}
\end{remark}

\noindent Both examples live \emph{outside} the three-manifold case of Proposition~\ref{prop:trivial-bundle-3-manifold}. Inside that case, the bundle is trivial and we lose no generality by treating the connection as a globally defined $A \in \Omega^1(M, \mathfrak{g})$, which is what Chapter~\ref{ch:chern-simons} does.

\begin{remark}[For readers who want the full bundle story]\label{rem:nakahara-pointer}
A connection on a principal bundle can be defined globally without choosing a section; the definition of the curvature is globally a horizontal equivariant $\mathfrak{g}$-valued two-form; the Bianchi identity is globally $D\Omega = 0$ for the bundle-theoretic covariant derivative $D$. The entire Chern--Simons machinery, in fact, has a globally-defined version that works on bundles that are \emph{not} trivial, at the cost of the action $S_{\CS}[A]$ being replaced by an object that depends on a choice of extension into a 4-manifold bounded by $M^3$. Dijkgraaf and Witten \cite{DijkgraafWitten1990} develop this global viewpoint for the purpose of classifying topological actions by group cohomology, a story we take up in Chapter~\ref{ch:dw}. For Chapter~\ref{ch:chern-simons}, the trivial-bundle hypothesis is enough.
\end{remark}


\section{Gauge transformations and the transformation of $F$}\label{sec:gauge-transf}

The connection one-form $A$ of Definition~\ref{def:connection-one-form} depends on the choice of local section $s : U \to P$ implicit in Definition~\ref{def:local-connection}. Changing $s$ is equivalent to acting by a smooth $G$-valued function $g : M \to G$, and the induced transformation of $A$ is what physicists call a \emph{gauge transformation}. This section derives the transformation law of $A$ and shows by direct computation that $F$ transforms by conjugation, $F^g = g^{-1} F g$, recovering the non-abelian analogue of the $U(1)$ gauge invariance of Example~\ref{ex:gauge-invariance-F}. The computation is the template for every gauge-transformation argument in Chapter~\ref{ch:chern-simons}, including the derivation of the Wess--Zumino piece in the transformation law of $\omega_{\CS}$; we execute it here with every sign tracked.

\subsection{Gauge transformation of the connection}

\begin{definition}[Gauge transformation]\label{def:gauge-transf}
Let $g : M \to G$ be a smooth map. The \emph{gauge transformation} of a connection $A \in \Omega^1(M, \mathfrak{g})$ by $g$ is
\begin{equation}\label{eq:gauge-A-2}
A \;\longmapsto\; A^g \;:=\; g^{-1} A g + g^{-1} dg.
\end{equation}
\end{definition}

The first term on the right of \eqref{eq:gauge-A-2} rotates $A$ by the adjoint action of $g^{-1}$; the second is an inhomogeneous piece that reflects the change of section. Together they make the parallel-transport operator $\mathcal{P} \exp \oint A$ transform by conjugation under gauge transformations, which is the invariance property physicists use in practice.

\begin{remark}[Notation for the Maurer--Cartan and conjugated connection]\label{rem:theta-tildeA}
For the rest of this section we write
\begin{equation}\label{eq:theta-tildeA-def}
\theta \;:=\; g^{-1} dg, \qquad \tilde A \;:=\; g^{-1} A g, \qquad A^g \;=\; \tilde A + \theta.
\end{equation}
The form $\theta$ is the \emph{Maurer--Cartan one-form} of the map $g : M \to G$; geometrically, it is the pullback to $M$ of the Maurer--Cartan one-form on $G$ along $g$. Chapter~\ref{ch:chern-simons} makes extensive use of its properties.
\end{remark}

\subsection{The identity $dg^{-1} = -g^{-1}(dg) g^{-1}$}

One lemma is used throughout.

\begin{lemma}\label{lem:dginv}
For any smooth $g : M \to G$,
\begin{equation}\label{eq:dginv-lemma}
dg^{-1} \;=\; -\, g^{-1}\, (dg)\, g^{-1}.
\end{equation}
\end{lemma}

\begin{proof}
Differentiate the identity $g^{-1} g = \mathbf{1}$ using the Leibniz rule for matrix products of $0$-forms:
\begin{equation*}
0 \;=\; d(g^{-1} g) \;=\; (dg^{-1})\, g \,+\, g^{-1}\, (dg).
\end{equation*}
Solve for $dg^{-1}$: move $g^{-1}(dg)$ to the left and then right-multiply by $g^{-1}$,
\begin{equation*}
(dg^{-1})\, g \;=\; -\, g^{-1}\, (dg) \quad\Longrightarrow\quad dg^{-1} \;=\; -\, g^{-1}\, (dg)\, g^{-1}. \qedhere
\end{equation*}
\end{proof}

\subsection{The Maurer--Cartan equation}

\begin{lemma}[Maurer--Cartan]\label{lem:maurer-cartan}
For any smooth $g : M \to G$, the Maurer--Cartan one-form $\theta = g^{-1} dg$ satisfies
\begin{equation}\label{eq:maurer-cartan}
d\theta \,+\, \theta \wedge \theta \;=\; 0.
\end{equation}
\end{lemma}

\begin{proof}
Apply $d$ to $\theta = g^{-1} dg$ using the graded Leibniz rule of Lemma~\ref{lem:ch2-leibniz-nilpotent}. Since $g^{-1}$ is a 0-form,
\begin{equation*}
d\theta \;=\; d(g^{-1} dg) \;=\; (dg^{-1}) \wedge (dg) \,+\, g^{-1}\, d^2 g \;=\; (dg^{-1}) \wedge (dg),
\end{equation*}
where $d^2 g = 0$ by the nilpotency identity \eqref{eq:ch2-nilpotent}. Substitute Lemma~\ref{lem:dginv}:
\begin{equation*}
(dg^{-1}) \wedge (dg) \;=\; \bigl(-\, g^{-1}\, (dg)\, g^{-1}\bigr) \wedge (dg).
\end{equation*}
In the matrix-valued-form notation of Section~\ref{sec:lie-forms}, the wedge is scalar on the form side and matrix-product on the Lie-algebra side; the factor $g^{-1}$ sitting between the two $dg$'s combines with the rightmost $(dg)$ to give $g^{-1}\, dg = \theta$, while the leftmost $g^{-1}\,(dg)$ is also $\theta$. So
\begin{equation*}
\bigl(-\, g^{-1}\, (dg)\, g^{-1}\bigr) \wedge (dg) \;=\; -\, \bigl(g^{-1}\, dg\bigr) \wedge \bigl(g^{-1}\, dg\bigr) \;=\; -\, \theta \wedge \theta,
\end{equation*}
proving $d\theta = -\theta \wedge \theta$, i.e.\ \eqref{eq:maurer-cartan}.
\end{proof}

\begin{remark}[The Maurer--Cartan equation as a statement about the group]\label{rem:mc-on-group}
The Maurer--Cartan equation is really a statement about the Lie group $G$ itself: the canonical left-invariant $\mathfrak{g}$-valued one-form $\theta_G \in \Omega^1(G, \mathfrak{g})$ on $G$ satisfies $d\theta_G + \theta_G \wedge \theta_G = 0$ as an identity of forms on $G$. Pulling back along any $g : M \to G$ gives \eqref{eq:maurer-cartan}. Chapter~\ref{ch:chern-simons} uses this form extensively: the Wess--Zumino three-form $\tr(\theta^{\wedge 3})$ that controls level quantization is the pullback along $g$ of a specific three-form on $G$, and \eqref{eq:maurer-cartan} is what makes that three-form closed.
\end{remark}

\subsection{The field strength transforms by conjugation}

We can now verify by direct calculation that $F^g = g^{-1} F g$. This is the non-abelian analogue of Example~\ref{ex:gauge-invariance-F} and is the template for the $\omega_{\CS}$ calculation in Chapter~\ref{ch:chern-simons}.

\begin{theorem}[Transformation of the field strength]\label{thm:F-transformation}
For any smooth $g : M \to G$ and any connection $A \in \Omega^1(M, \mathfrak{g})$, the field strength $F = dA + A\wedge A$ transforms as
\begin{equation}\label{eq:F-transformation}
F^g \;=\; g^{-1}\, F\, g.
\end{equation}
\end{theorem}

\begin{proof}[Sketch]
Expand $F^g = dA^g + A^g\wedge A^g$ using $A^g = \tilde A + \theta$. The exterior derivative produces $d\tilde A - \theta\wedge\theta$ after invoking Maurer--Cartan, and the square expands to $\tilde A\wedge\tilde A + \tilde A\wedge\theta + \theta\wedge\tilde A + \theta\wedge\theta$ with $\tilde A\wedge\tilde A = g^{-1}(A\wedge A)g$. In the sum, three pairs of mixed terms cancel: the two copies of $\theta\wedge\theta$; the $\pm\theta\wedge\tilde A$ pair; and the $\pm\tilde A\wedge\theta$ pair. What remains is $g^{-1}(dA + A\wedge A)g = g^{-1}Fg$. The same cancellation pattern, with a different coefficient on the cubic piece, drives the $\omega_{\CS}$ gauge-transformation calculation of Chapter~\ref{ch:chern-simons}, and the fully explicit version of that argument lives there.
\end{proof}

\begin{remark}[Infinitesimal gauge transformation]\label{rem:infinitesimal-gauge}
For a gauge transformation close to the identity, $g = \mathbf{1} + \varepsilon$ with $\varepsilon = \varepsilon^a t^a$ an infinitesimal Lie-algebra-valued 0-form, \eqref{eq:gauge-A-2} gives, to first order in $\varepsilon$,
\begin{equation*}
A^g \;=\; (\mathbf{1} - \varepsilon) A (\mathbf{1} + \varepsilon) \,+\, (\mathbf{1} - \varepsilon)\, d\varepsilon \;=\; A \,+\, [A, \varepsilon] \,+\, d\varepsilon \,+\, O(\varepsilon^2).
\end{equation*}
The infinitesimal variation is $\delta A = d\varepsilon + [A, \varepsilon]$, which is the \emph{covariant derivative} of a Lie-algebra-valued 0-form with connection $A$. Differentiating $F^g = g^{-1}Fg$ at $g = \mathbf{1}$ correspondingly gives $\delta F = [F, \varepsilon]$. This infinitesimal form appears implicitly in Schwartz's non-abelian-gauge-theory chapter; the finite, globally-defined version \eqref{eq:F-transformation} is needed when we integrate over large gauge transformations in Chapter~\ref{ch:chern-simons}.
\end{remark}

\medskip

\noindent This completes Chapter~\ref{ch:forms}. The differential-forms toolkit of Section~\ref{sec:forms-basics}, the integration and Stokes results of Section~\ref{sec:integration-stokes}, the Lie-algebra-valued extension of Section~\ref{sec:lie-forms}, the principal-bundle vocabulary of Section~\ref{sec:principal-bundles}, and the gauge-transformation law of this section together form the geometric language used in Chapters~\ref{ch:chern-simons}, \ref{ch:bf}, \ref{ch:dw}, and \ref{ch:gensym}. Chapter~\ref{ch:axiomatic} pivots back to the topological side: the Atiyah--Segal axiomatic framework of TQFT builds on top of the same orientation, integration, and cohomological structure we have just developed.
